% Prose of the robustness follow-ups RF1-RF3 (after the second review). Numbers are macros
% from scripts/report_robust.py; qualitative statements are checked by scripts/check_claims.py.

\newcommand{\RobustIntro}{%
A second review raised two methodological questions about study~2---whether the
response-derived clusters favour EOST-M, and whether the Bayesian CATs were estimated with
enough Markov chain Monte Carlo---and asked whether the effects generalize across benchmarks
under a less conservative criterion than the frozen one. Three follow-ups answer them (RF1--RF3);
their plans and decision rules were written in the research log before they were run, and
none of them changes a frozen decision. Further, exploratory follow-ups of the MCMC question
(RF2b--RF2d) were designed after the diagnostics of RF2 had been seen, and a sensitivity
analysis of RF3 (RF3b) after a third review; their plans, too, were written before they were
run.}

\newcommand{\RobustClusterText}{%
The strict clusters join systems
whose 2PL residuals correlate, and those residuals are computed from every system's full
response vector, the target's included. Excluding the calibrators that resemble a target
could make its panel harder to calibrate and so favour a method built for calibration
uncertainty. We kept the frozen targets, holdouts and panel seeds and changed only the systems
excluded from each panel: heuristic families alone (no response data), leave-bank-out
clusters (residual correlations computed on the other benchmarks that contain the same
systems), and the frozen construction at thresholds 0.4, 0.6 and 0.7. Recomputing the frozen
construction reproduced all \RobReproRows\ frozen rows exactly. Under every alternative rule,
R2 and C1 kept their frozen median improvements (\RobRTwoMedMin\% and \RobCOneMedMin\%), R2
remained positive on all eight benchmarks with \RobRTwoFavRange\ of them Holm-significant, and
C1 remained positive on seven with \RobCOneFavRange\ significant, none
against (Table~\ref{tab:robclust}; per benchmark in Table~\ref{tab:robclustbanks}): the
robustness criterion fixed in advance was \RobRuleMet. The response-derived links matter
only on the Open LLM Leaderboard banks, where hundreds of fine-tunes share a few base models;
on the other benchmarks the clusters are essentially the heuristic families
(Table~\ref{tab:robclustcount}). Adaptivity (C4) had \RobCFourFavRange\
significant benchmarks under the alternative rules (frozen: \STwoCFourFav), with medians between
\RobCFourMedMin\% and \RobCFourMedMax\%.}

\newcommand{\RobustMetaText}{%
The frozen rule asks for significance on at least
five of eight benchmarks after a Holm correction over 88 tests, so an effect can favour EOST-M
on all eight benchmarks and still be declared not supported. A random-effects meta-analysis
asks a different, secondary question---does an effect point the same way across
benchmarks?---without changing any decision (Table~\ref{tab:meta}). Treating the benchmarks
as independent, the pooled improvement excluded zero for \MetaNExclZero\ of the 11
superiority hypotheses, including the \MetaNExclZeroUnsup\ that missed the frozen threshold by
one benchmark (\MetaUnsupList; $I^2$ \MetaUnsupIMin--\MetaUnsupIMax\%), and the 95\%
prediction interval for a new benchmark excluded zero only for \MetaPiExclList. But five
benchmarks share the same 395 Open LLM Leaderboard systems, a dependence the model ignores.
Counting them as one source (RF3b: \MetaSrcK\ sources, standard error from a bootstrap over the
shared systems), \MetaSrcNExclZero\ of the 11 intervals excluded zero (\MetaSrcExclList);
\MetaSrcLostList\ stayed positive in \MetaSrcLostPosMin--\MetaSrcLostPosMax\ of the
\MetaSrcK\ sources but with intervals including zero, so of the four near misses only
\MetaSrcUnsupKeptList\ kept an interval above zero. The dependence through shared systems was
small (median correlation \MetaSrcR); the loss comes from having four sources and from their
heterogeneity. Most effects thus point the same way across benchmarks; the analysis does not
show that they generalize to a new one.}

\newcommand{\RobustMCMCDesign}{%
On six study-2 targets per benchmark, drawn with a fixed seed (replicate~0, $N=8$;
\ConvPanels\ panels), we ran the post hoc copy of the frozen sampler, which reproduces it
exactly, with four chains. Chain~1 started at EOST-M's estimates, as the frozen arm did;
chains 2--4 started from dispersed points (abilities perturbed by $\mathcal N(0,0.5^2)$;
$\mu_b\pm0.5$; $\sigma_b\times0.7$ or $1.4$; $\mu_a\pm0.3$; $\sigma_a\times0.7$ or $1.4$).
Each chain ran 250 burn-in sweeps and 300 further sweeps, every third of which was retained
(400 retained draws in all). Diagnostics use the 300 post-burn-in sweeps of each chain:
rank-normalized split $\hat R$ (the larger of the bulk and folded versions), bulk effective
sample sizes (ESS) with Geyer's initial monotone sequence, and Monte Carlo standard errors
from the ESS of the raw draws \citep{vehtari2021rank,geyer1992practical}. The posterior-draw
marginal arms average the item posteriors over the 400 pooled draws; the joint arms mix over
every tenth draw (40). Table~\ref{tab:rf2} gives the convergence per benchmark,
Table~\ref{tab:rf2arms} the arms on the same cells and Figure~\ref{fig:rf2traces} the traces
of the first panel of each benchmark.}
\newcommand{\RobustMCMCFrozen}{%
The overall location converged ($\hat R<1.05$ on \ConvLocRhatFive\% of panels, median bulk
ESS \ConvLocEssMed), and so did $\mu_b$ on most panels (\ConvMuBRhatFive\%), but $\sigma_b$,
$\mu_a$ and $\sigma_a$ did not: $\hat R<1.05$ on at most \ConvHypRhatFiveMax\% of panels,
median $\hat R$ \ConvHypRhatMedMin--\ConvHypRhatMedMax, median bulk ESS
\ConvHypEssMedMin--\ConvHypEssMedMax, and Monte Carlo standard errors of
\ConvHypMcseMin--\ConvHypMcseMax\% of their posterior standard deviations. Neither did the
ability scale, the standard deviation of the calibrators' abilities ($\hat R<1.05$ on
\ConvScaleRhatFive\% of panels): the ratio of the largest to the smallest chain mean had a
median of \ConvScaleRatioMed\ (at most \ConvScaleRatioMax). Multiplying all abilities and
difficulties by a constant and dividing the discriminations by it leaves the likelihood
unchanged, so with eight calibrators only their $\mathcal N(0,1)$ prior fixes the scale of
the calibration; the frozen sampler has a translation move but no move along this ridge, on
which the spread of the calibrators' abilities moved together with $\sigma_b$ (median
correlation \ConvCorScaleSdB) and $\mu_a$ (\ConvCorScaleMuA). The four chains' held-out log
losses differed by more than 5\% on \ConvLLSpreadGtFive\% of panels (at most
\ConvLLSpreadMax\%). Under the rule fixed in advance the study-2 statement about the
Bayesian CATs was \ConvRule: on individual cells the long chains changed the held-out log
loss by a median of \ConvLongAbsVsSixtyMin--\ConvLongAbsVsSixtyMax\% relative to a burn-in of
60, more than the 0.5\% allowed, although the signed median changes were small
(\ConvLongVsSixtyMin\ to \ConvLongVsSixtyMax\%). As the better estimate, the long-chain arms
differed from EOST-M by a median of \ConvLongVsEostMin\ to \ConvLongVsEostMax\% in held-out
log loss, with the benchmarks split (Table~\ref{tab:rf2arms}); on these cells they covered
the score \ConvLongCovRange\ of the time, against \ConvEostCov\ for EOST-M.}
\newcommand{\RobustMCMCScale}{%
Exploratory follow-ups, designed after these diagnostics and each planned before it was run,
added the missing moves. RF2b added to every sweep a Metropolis--Hastings move along the
scale ridge, applied to the posterior of the abilities and hyperparameters with the item
parameters summed out on the grid: the abilities, $\mu_b$ and $\sigma_b$ are multiplied by
$s$ and $\mu_a$ is decreased by $\log s$, with $\log s=\pm0.2$ or $\pm0.4$ (multiples of the
step of the log-discrimination grid, so that the move maps that grid onto itself); the
acceptance ratio contains the Jacobian $s^{N+2}$, and the item parameters are then drawn from
their exact conditional, which makes the scheme a partially collapsed Gibbs sampler
\citep{vandyk2008partially} with a discrete form of parameter expansion
\citep{liu1999parameter}. With this move (median acceptance \SclAccMed\%; Table~\ref{tab:rf2x}),
$\hat R<1.05$ was reached on \SclScaleRhatFive\% of panels for the ability scale, \SclSdBRhatFive\% for
$\sigma_b$ and \SclMuARhatFive\% for $\mu_a$ (RF2: \ConvScaleRhatFive, \ConvSdBRhatFive\ and
\ConvMuARhatFive\%), but only on \SclSdARhatFive\% for $\sigma_a$ (median bulk ESS
\SclSdAEssMed), so the convergence criterion fixed for RF2b was not met; the four chains'
held-out log losses now differed by at most \SclLLSpreadMax\% (RF2: \ConvLLSpreadMax\%). Its
arms differed from EOST-M by a median of \SclVsEostMin\ to \SclVsEostMax\% in held-out log
loss (descriptive, since the chains had not converged).
\par
The frozen sampler's conjugate hyperparameter draws, which RF2b kept, mix slowly when each
item carries little information, and they treat the grid prior as a continuous density: they
condition on the sampled grid points, whereas EOST-M normalizes the prior over the grid, and
the two differ when a standard deviation approaches the grid step (0.2 for $\log a$), where the unnormalized version accumulates at the lower bound of $\sigma_a$
(\ConvFloorChains\ of RF2's \ConvChains\ chains and \SclFloorChains\ of RF2b's \SclChains\
spent most of their sweeps there; Figure~\ref{fig:rf2traces}, GSM8K). RF2c and RF2d
therefore sample the model of Lemma~\ref{lem:exact} exactly. RF2c replaces the conjugate
draws by random-walk Metropolis moves on $\mu_a$, $\log\sigma_a$, $\mu_b$ and $\log\sigma_b$
with the item parameters summed out (two rounds per sweep; cheap because the grid is the
product of a discrimination axis and a difficulty axis; step sizes adapted during burn-in
only), keeps the scale move and the frozen ability update, and multiplies the weights of the
translation move by the factor $(Z(\eta)/Z(\eta'))^I$ that the normalized prior requires.
RF2c was stopped after \RfcPanels\ panels: on GSM8K every recorded quantity converged
(largest $\hat R$ \RfcGsmRhatMax), but on ARC-C and HellaSwag the overall location did not
($\hat R$ up to \RfcLocRhatMax), because when the difficulty prior extends beyond the edge of
the grid (HellaSwag) the corrected translation weights, raised to the power of the number of
items, block almost every shift. RF2d therefore adds a collapsed translation---the abilities
and $\mu_b$ shifted by 1, 2 or 4 grid steps with the item parameters summed out, which items
at the edge of the grid cannot block---and finer scale steps
($\log s\in\{\pm0.05,\pm0.1,\pm0.2,\pm0.4\}$): it moves the hyperparameters, the scale and
the location by Metropolis--Hastings steps with the item parameters summed out. All samplers
were checked before they were run: with the new move switched off, RF2b's sampler reproduces
RF2's draws exactly, and the stationary laws of the scale moves, the hyperparameter moves and
both translation moves match exact computations on small problems. RF2b used a burn-in of
200 sweeps and RF2d one of 150, each followed by 200 sweeps per chain with every second
retained (400 draws).}
\newcommand{\RobustMCMCExact}{%
RF2d came much closer to convergence (median acceptance: translation \ExtAccTrans\%, scale
\ExtAccMed\%, hyperparameter moves about \ExtAccMuA\%). $\hat R<1.05$ held on
\ExtLocRhatFive\% of panels for the location, \ExtScaleRhatFive\% for the ability scale, and
\ExtMuBRhatFive\%, \ExtSdBRhatFive\%, \ExtMuARhatFive\% and \ExtSdARhatFive\% for $\mu_b$,
$\sigma_b$, $\mu_a$ and $\sigma_a$, with medians of $\hat R$ between \ExtLocRhatMed\ and
\ExtRhatMedMaxAll\ across these quantities, a largest value of \ExtRhatMaxAll\ (against
\ConvHypRhatMax\ in RF2), and median bulk ESS between \ExtEssMedMin\ and \ExtLocEssMed\
(Table~\ref{tab:rf2x}, Figure~\ref{fig:rf2dtraces}). No chain stayed at the lower bound of
$\sigma_a$, but on \ExtBoxSdB\ panels $\sigma_b$ and on \ExtBoxSdA\ panels $\sigma_a$ spent
more than a quarter of the sweeps within 5\% of the upper bound of the prior box inherited
from the frozen sampler (4 and 1.5), which therefore constrains the posterior there. With
$\hat R<1.05$ on only \ExtRhatFiveMin\% of panels for the slowest quantity, short of the 90\%
fixed in advance, the convergence criterion was \ExtConverged, and a fully converged
comparator was not obtained. The four chains' held-out log losses differed by a median of
\ExtLLSpreadMed\% (at most \ExtLLSpreadMax\%). Under the model of Lemma~\ref{lem:exact}, the
posterior mean of the ability scale exceeded EOST-M's point estimate on \ExtPostAboveEb\% of
panels (median ratio \ExtPostOverEbMed): the calibration that EOST-M conditions on has a
smaller spread than the posterior.}
\newcommand{\RobustMCMCConclusion}{%
Across all \AllBayesNArms\ combinations of Bayesian arm and sampler on these cells, from the
frozen chain to RF2d, the median difference from EOST-M in held-out log loss stayed within
$\pm$\AllBayesVsEostMaxAbs\%, with the benchmarks split (RF2d: \ExtVsEostMin\ to
\ExtVsEostMax\%; Table~\ref{tab:rf2arms} gives every arm). The study-2 comparison did not
move as the sampler improved, which makes an artefact of the frozen chain unlikely; our
statements about the Bayesian CATs nevertheless concern the implementations studied, not an
exactly computed fully Bayesian test. The RF2d arms covered the score as often as EOST-M did
on the same cells (\ExtCovMin--\ExtCovMax\ against \ExtEostCov), so integrating the
calibrators' abilities and the hyperprior did not repair the under-coverage; the higher
coverage of RF2's long chains (\ConvLongCovMin) did not survive the exact model. In score
error the RF2d arms were lower than EOST-M by a median of \ExtIaeGainMin--\ExtIaeGainMax\%
over cells but on only \ExtIaeBanksMin--\ExtIaeBanksMax\ of 8 benchmarks, a possible
advantage that 48 cells cannot establish.}
